%%
%% (Somewhat) Practical Anonymous Router
%%
\section{(Somewhat) Practical Anonymous Router}
\label{sec:spar}
\gls{niar} introduced the anonymous router, a primitive that scales better than existing solutions with the same privacy guarantees~\cite{cryptoeprint:2021/435}. The router construction itself makes the communication cost practical. Each client contacts only the router and the router contacts each client, so one round costs $\eta$ connections. Its instantiation, however, is hard, if not impossible, to realize in practice (\autoref{chp:background}).

The \gls{spar} makes the instantiation practical by building the router on multi-party \gls{fhe}~\cite{cryptoeprint:2025/860}. The server permutes the messages by computing on their ciphertexts, so the permutation exists only under encryption. It is derived from random values the clients generate locally, fresh in every round; the setup takes no permutation as input and involves no trusted party. Security rests on the \gls{rlwe} assumption of \autoref{sec:fhe}, for which practical implementations already exist. The price is a second client-server round trip; no party holds the joint secret key, so decryption requires every client to partially decrypt the results (\autoref{sec:mpfhe}).

A final measure of practicality remains; how long does it take to run the protocol? This thesis follows the initial preprint of May 2025~\cite{cryptoeprint:2025/860}, the only version available at the time of writing, which has no associated implementation.

%%
%% SYSTEM MODEL
\subsection{System Model}
\label{sec:spar:model}
The system contains $\eta$ clients and a single server, with communication happening in \emph{rounds}. Every client submits exactly one message per round, sending a random message if it has nothing to say, and all traffic passes through the server. The server is untrusted for privacy but trusted for availability, and may collude with any number of corrupted clients. Sender anonymity follows from the IND-CPA security of the underlying scheme, with definitions and proofs given in the original paper~\cite{cryptoeprint:2025/860}.

Messages are assumed to be uniformly distributed, so that placement at a random encrypted index stands in for a shuffle. Messages sharing a bucket must additionally be sorted obliviously before the list is opened, since their relative order would otherwise reveal which sender wrote first.

%%
%% THE PROTOCOL
\subsection{Protocol}
\label{sec:spar:protocol}
A one-time setup is followed by polynomially many rounds under the same keys, each consisting of four computations separated by two client-server round trips.

\subsubsection{Setup Phase}
The clients agree on a set of \gls{fhe} parameters and run the multi-party key generation of \autoref{sec:mpfhe}, producing a joint public key under which every round is encrypted.

\subsubsection{Online Phase}
\begin{itemize}
    \item \emph{Client: Encrypt} Client $i$ encrypts its message under the joint key together with three bucket indices $a_0, a_1, a_2 \overset{\$}{\leftarrow} [0, \eta)$ sampled uniformly. A single index per client would collide too often, whereas with $\eta$ buckets of three slots the probability that all nine candidate slots are taken falls below $10^{-11}$, and on that event the message is dropped and retransmitted in a later round~\cite{cryptoeprint:2025/860}.
    \item \emph{Server: Write} The server keeps a message matrix $L$ of $\eta\times3$ \gls{rlwe} slots and an availability matrix $I$ of $\eta\times3$ encrypted flags marking which slots are free. Each index enters the write as an indicator vector of $\eta$ encrypted bits where exactly one encrypts $1$ (\autoref{sec:spar:bandwidth}). For each client, the write loop of \autoref{alg:spar:homplacing} places the message in the first free slot among the three chosen buckets, tracked by the encrypted flag $\mathsf{hasWritten}$. Once all $\eta$ messages are written, the slots of each bucket are sorted obliviously and $L$ is broadcast to the clients.
    \item \emph{Client: Decrypt} Each client partially decrypts all $3\eta$ slots with its own share (\autoref{sec:mpfhe}) and sends the result back to the server.
    \item \emph{Server: Decrypt} The server aggregates the partial decryptions of each slot (\autoref{sec:mpfhe}) and recovers the list. Empty slots decrypt to zero and are discarded, leaving the $\eta$ messages in an order that is a fresh uniform permutation of the senders.
\end{itemize}

%%
%% BANDWIDTH
\subsection{Bandwidth Optimizations}
\label{sec:spar:bandwidth}
Transmitted directly, the three indicator vectors of one client occupy $3\eta$ ciphertexts against a single one for its message. The paper instead uploads each index as its $\log\eta$ bits, encrypted separately, and lines 1--12 of \autoref{alg:spar:homplacing} rebuild the vectors on the server. A binary tree of depth $\log\eta$ starts from a root encrypting $1$, each level multiplies the current node by the index bit and by its complement, and the $\eta$ leaves hold the indicator vector. The write loop from line 13 onward is unchanged whether the vectors are received or rebuilt.

The products consume the index bits in \gls{rgsw} form, and an \gls{rgsw} ciphertext consists of $2\ell$ \gls{rlwe} rows \eqref{def:rgsw}. Clients therefore upload plain \gls{rlwe} encryptions, which the server expands with $\mathsf{HomExpand}$~\cite{chillotti:onionring:2019, park:sortinghat:2022, cryptoeprint:2025/1088}. When $N > 3\eta$, the $3\eta$ partial decryptions of one client pack into a single ciphertext through the ciphertext conversion of~\cite{cryptoeprint:2020/015}.

%%%
%%% sPAR Algorithm 2 - HomPlacing* for the choice of three slots
%%% Paper-faithful (eprint 2025/860); generic homomorphic products written
%%% "\cdot" as in the paper --- the typed instantiation is thesis material (Ch. 3).
%%%
\begin{algorithm}[ht!]
\caption{Homomorphic placing $\mathsf{HomPlacing}^*$ for the choice of three slots~\cite{cryptoeprint:2025/860}}
\label{alg:spar:homplacing}
\begin{algorithmic}[1]
    \Require Encrypted value $V_r$; encrypted indices $A_0, A_1, A_2$, where $A_d$ is $\log\eta$ encrypted bits $\{c_{d,0}, \dots, c_{d,\log\eta - 1}\}$
    \Require Server state matrices $L_{\eta\times3}$ and $I_{\eta\times3}$
    \Ensure Encrypted boolean indicating whether the value was placed
    \State $z_{i,d} := 0$ for $i \in [0, \eta),\; d \in \{0, 1, 2\}$
    \For{$d \gets \{0, 1, 2\}$}
        \State $b_0 := 1$ and $b_i := 0$ for $i \in [1, 2\eta - 1)$
        \Comment{Empty binary tree with $\eta$ leaves}
        \For{$i \gets [0, \log\eta)$}
            \For{$j \gets [0, 2^i)$}
                \State $r := b_{2^i - 1 + j}$
                \State $b_{2^{i+1} - 1 + 2j} := r - r\cdot c_{d,i}$
                \State $b_{2^{i+1} + 2j} := r\cdot c_{d,i}$
            \EndFor
        \EndFor
        \State $z_{i,d} := b_{\eta - 1 + i}$ for $i \in [0, \eta)$
        \Comment{Leaves: indicator vector of $a_d$}
    \EndFor
    \State $\mathsf{hasWritten} := 0$
    \For{$d \gets \{0, 1, 2\}$}
        \For{$k \gets \{0, 1, 2\}$}
            \Comment{Slots per bucket}
            \For{$i \gets [0, \eta)$}
                \State $h := z_{i,d}\cdot I_{i,k}\cdot(1 - \mathsf{hasWritten})$
                \State $L_{i,k} := L_{i,k} + h\cdot V_r$
                \State $I_{i,k} := I_{i,k} - h$
                \Comment{Mark the slot as used}
                \State $\mathsf{hasWritten} := \mathsf{hasWritten} + h$
            \EndFor
        \EndFor
    \EndFor
    \State\Return $\mathsf{hasWritten}$
\end{algorithmic}
\end{algorithm}
 % paper Alg. 2: choice-of-three (new float)
\newpage

%%
%% EXTERNAL PRODUCTS
\subsection{Depth Requirements}
\label{sec:spar:requirements}
\autoref{alg:spar:homplacing} is stated for a generic homomorphic multiplication, but the circuit it describes is deep. Every level of the binary tree multiplies the running node by a further encrypted bit, and the write loop chains $9\eta$ multiplications through $\mathsf{hasWritten}$ for each message. Ciphertext multiplication in the sense of \autoref{sec:bgv} would exhaust the noise budget long before a round completed, so the authors propose evaluating the tree with the external product, whose noise grows additively along such a chain (\autoref{rem:extprod-asymmetry}).
